\subsubsection{货箱移除、重新插入与局部调整}
本文采用自适应大邻域搜索的破坏--修复框架\cite{ropke2006alns}，把标准框架中的操作对象具体化为货箱、访问顺序和飞机--电池任务链。一次破坏使用随机移除3--7箱、移除一至两项完整任务、移除空间相关货箱，或从当前平均负载最重机型的长任务中移除货箱。删除服务区后也重新计算航段；山区几何不保证直接删点后的航段代价更小，不能跳过物理核验。

修复时，对每个待插入货箱尝试现有任务及单独新建任务，允许重新选择A/B/C机型。单个新增服务区遍历原访问顺序中的插入位置；需要重构多个服务区时，至多5个不同服务区枚举排列，更多服务区仅保留最近邻顺序及反序。后一规则是搜索筛选，不是题目对访问数量的限制。先按平均机型工作量、能耗和时限必要条件排序，保留前10项及每种机型最好的1项，再对去重后的候选遍历该机型任务链中的插入位置，执行式\eqref{eq:q2_decoder}的完整排程。因而这里没有声称穷举全部插入组合。

设货箱$b$的可行插入评分从小到大为$F_{b1},F_{b2},\ldots$。贪心修复选择当前最低评分；$k$阶后悔修复优先处理
\begin{equation}
 R_k(b)=\sum_{j=2}^{k}(F_{bj}-F_{b1}),\qquad k\in\{2,3\}
 \label{eq:q2_regret}
\end{equation}
较大的货箱，即优先安排缺少好替代位置的货箱。可用位置不足$k$个时重复最后一项评分。每轮修复后再进行20次局部提议：同机型任务换序、任务改型、跨任务换箱和单箱迁移的选择概率依次为0.40、0.18、0.22、0.20；全部提议仍需重算物理属性和真实资源日程。

\textbf{物理等价货箱的原始编号匹配。}同区同类箱子可能质量、体积相同而截止要求不同，求解器不能只关心总质量。实现先依据当前日程生成逐箱交付事件，再在同一服务区、同一物资类型的组内，把原始箱号按“硬截止、期望送达、箱号”排序，把交付事件按实际时刻排序，依次匹配。无硬截止记为排序中的正无穷。当前输入中各匹配组的质量和体积一致，故交换身份不改变任务箱数、载荷、装卸时间和能耗；每箱的原始时限及权重始终随原始箱号保留。匹配后重新排程并逐箱验收，而不直接宣称该排序解决了任意加权迟到问题。若输入不满足组内物理等价，则必须细分组，不能沿用此交换。

\subsubsection{搜索接受规则与独立方案核验}
最终方案始终按式\eqref{eq:q2_objective}比较；为在搜索过程中探索不同邻域，程序另采用标量代理评分。记$n_H$为硬迟到箱数、$L_H$为硬迟到秒数之和，则工期搜索模式下
\begin{equation}
\begin{split}
 F(S)={}&2\times10^6n_H+2000L_H+100L_w+T_{\max}\\
        &+1.5E+2N+2\times10^{-4}\sum_bC_b .
\end{split}
\label{eq:q2_search_score}
\end{equation}
式中时间用秒、能耗用kWh；系数是代码固定的搜索导航设置，不是将不同量纲变为等价救援成本。硬迟到候选即使参与探索，也不能进入正式可行方案档案。

令迭代上限为$K$、当前轮为$k$。大邻域和局部提议分别使用温度
\begin{equation}
 \Theta_k=55(1-k/K)^2+1,\qquad
 \theta_k=25(1-k/K)^2+0.2.
 \label{eq:q2_temperature}
\end{equation}
评分更低的候选直接接受，否则以$\exp[-(F(S')-F(S))/\Theta_k]$的概率接受；局部操作将$\Theta_k$替换为$\theta_k$。这使当前搜索状态能够暂时变差，但最好可行方案仅在词典序严格改善时更新。

破坏与修复算子的初始权重为1，按权重比例抽取。某大邻域操作得到新的最好可行解、只改善当前代理评分、或仅被接受时，奖励依次为20、6、1；未接受为0。每25轮对使用过的算子$o$更新
\begin{equation}
 w_o\leftarrow\max\{0.15,\;0.8w_o+0.2\bar r_o\},
 \label{eq:q2_operator_weight}
\end{equation}
其中$\bar r_o$是本周期每次使用的平均奖励；未使用时保持其权重。每80轮把当前状态重置为最好可行解，属于搜索回退而非从零重启。达到迭代上限或墙钟上限之一即停止。

\begin{table}[htbp]\centering\small
\caption{任务重构与双资源搜索的一次完整迭代}\label{tab:q2_algorithm}
\begin{tabularx}{\linewidth}{cX}\toprule
步骤 & 输入、操作及更新对象\\\midrule
1 & 读取完整任务序列，重算物理属性、原始箱号交付及飞机--电池日程；保存当前状态与最好可行方案。\\
2 & 依据权重选择破坏、修复算子，移除相应货箱；按候选筛选规则生成插入位置、机型与访问顺序。\\
3 & 用贪心或后悔规则修复，匹配物理等价原始箱号，并完整核验容量、能量、时限和资源接续。\\
4 & 按代理评分和降温规则更新当前状态；仅对可行且词典序更好的解更新最优档案并记改进日志。\\
5 & 提出20次局部调整，再次核验与接受；按25轮和80轮周期分别更新算子权重和回退当前状态。\\
6 & 达到轮数或时间限制后，输出最好可行方案；独立程序从原始参数重新核验，而不复用主评价器的中间结果。\\\bottomrule
\end{tabularx}
\end{table}

\textbf{正式方案与新增试验的来历。}已保存的输入基线含23架次，工期5693.231489\,s、能耗66.225670\,kWh；选定记录标记为种子9201的后续候选。核对前后决策，可识别S001两项C型任务的一只食品箱与一只卫生用品箱交换，净节能0.011211\,kWh，工期不变。当前归档可以逐项重放这两份方案，但未保存这次历史搜索的完整逐轮日志和实际停止轮数，不能据此编造从空任务集开始的收敛曲线。后文20次试验采用另一共同受扰起点和150轮预算，1\%、2\%能耗备选各250轮；这些实际设置与历史最好方案明确分开。标准方法、候选来源和复现范围见附录\ref{app:source_provenance}。

独立验证逐项重建节点、货箱、航段、能耗、充电和交付，并检查每个飞机和电池区间。正文表\ref{tab:q2_schedule}、箱号附表与甘特图均由同一份23架次、80箱交付记录生成，原始决策列表顺序不再与按执行时刻排列的编号混用。
